\section{Теорема Эберлейна-Шмульяна}\label{seq:eberlein-smulyan}
\begin{theorem}[Эберлейн, Шмульян]
    Пусть $(X, \|\cdot\|)$ -- нормированное пространство и $K \subset X$ -- слабый компакт.
    Тогда $K$ -- слабый секвенциальный компакт.
\end{theorem}
\begin{proof}
    Для начала покажем, что $K$ является ограниченным по норме.
    Пусть $M = \sup_{x \in K} \|x\|$.
    Очевидно, что существует максимизирующая последовательность $\{x_n\} \subset K$ такая, что
    \[
        \lim_{n \rightarrow \infty} \|x_n\| = M.
    \]
    В тоже время, так как в слабой топологии $f \in X^*$ непрерывны (просто по её определению), то $f(K) \subset \Cm$ является компактом.
    А значит $f(K)$ ограничено и, следовательно, лежит в круге некоторого радиуса $M_f > 0$ с центром в нуле.
    Так как это верно для всех точек $K$, то верно и для максимизирующей последовательности, то есть
    \[
        \forall f \in X^* \hookrightarrow |f(x_n)| \leq M_f.
    \]
    Следовательно, последовательность $\|x_n\|$ является равномерно ограниченной (по теореме Банаха-Штейнгауза) и $M < +\infty$.

    Рассмотрим произвольную $\{x_n\} \subset K$.
    Пусть $L = [\Lin\{x_n\}_{n = 1}^{\infty}]_{\|\cdot\|}$.
    Тогда $L$ -- замкнутое по норме подпространство $X$ и $\{x_n\} \subset L \cap K$.
    Очевидно, что $L$ -- сепарабельно, мы можем взять множество вида
    \[
        M = \left\{\sum_{n = 1}^{N} \alpha_n x_n \mid N \in \N, \alpha_i \in \Q + i\Q\right\}.
    \]
    В то же время, $L$ является и слабо замкнутым подпространством $X$, поскольку $\forall x \in X \setminus L$
    по следствию из теоремы Хана-Банаха существует непрерывный функционал $g$ такой, что $g|_L = 0$ и $g(x) = 1$, а его норма $\|g\| = \frac{1}{\mathrm{dist}(x, L)}$.
    И тогда точка $x$ отделяется от $L$ окрестностью $V(x, g, 1)$, поскольку
    \[
        \forall z \in V(x, g, 1) \hookrightarrow |g(x)| - |g(z)| \leq |g(z) - g(x)| < 1 \Rightarrow |g(z)| > 0.
    \]
    Поэтому $X \setminus L$ является слабо открытым, а $L$, соответственно, слабо замкнутым.

    Заметим, что если у нас существует подпоследовательность $\{x_{n_k}\}$, которая слабо сходится к некоторому $x_0 \in X$, то $x_0 \in K$ в силу слабой замкнутости $K$.
    Соответственно, мы можем работать только с функционалами из $L^*$.

    Поскольку $L$ -- сепарабельно, то *-слабая топология в $L^*$ метризуема на шаре.
    С другой стороны, по теореме Банаха-Алаоглу, $B_1^*(0)$ является *-слабо компактным.
    А значит он является метрическим компактом, и, следовательно, сепарабельным.
    Пусть $\{f_m\}_{m = 1}^{\infty} \subset B_1^*(0)$ -- счётное всюду *-слабо плотное множество в шаре.

    Приступим к построению искомой подпоследовательности.
    Для всякого $x_n\colon$
    \[
        |f_1(x_n)| \leq \|f_1\|\|x_n\| \leq 1 \cdot M.
    \]
    Поэтому $\{f_1(x_n)\}$ -- ограниченная последовательность, поэтому допускает выделение сходящейся подпоследовательности $x_{n_{k_1}}$.
    Аналогично, последовательность $\{f_2(x_{n_{k_1}})\}$ будет ограничена, и мы снова сможем выделить сходящуюся подпоследовательность.
    Продолжая диагональный процесс Кантора мы получим $\{x_{n_m}\}$ такую, что $\forall f_k$ у нас имеется сходимость $f_k(x_{n_m})$ (куда-то).
    Если $\{x_{n_m}\}$ содержит конечное множество различных точек, то у нас есть стационарная подпоследовательность.
    Пусть это не так.
    Тогда, поскольку все точки принадлежат $K \cap L$ -- слабому компакту, а он является компактом предельной точки (то есть всякое бесконечное множество точек из $K \cap L$ имеет в нём предельную точку), то $\exists x_0 \in K \cap L$ -- предельная точка $E = \{x_{n_m}\}$.
    Заметим, что раз для всякого $f_k$ верно, что $f_k(x_{n_m})$ куда-то сходится, то оно сходится к $f_k(x_0)$.
    Покажем это.
    В силу фундаментальности $f_k(x_{n_m}):$
    $\forall \epsilon > 0 \forall k \in \N \exists m(\epsilon, k)$ такое, что $\forall r, s \geq m(\epsilon, k) \hookrightarrow$
    \[
        |f_k(x_{n_s}) - f_k(x_{n_r})| \leq \epsilon.
    \]
    Рассмотрим $U(x_0) = V(x_0, f_k, \epsilon) \setminus \{x_{n_m} \mid m = \overline{1, m(\epsilon, k)}, \ x_{n_m} \neq x_0\}$.
    Тогда существует $m_0 \in \N$ такое, что $x_{n_{m_0}} \in U(x_0)$ (в силу предельности точки $x_0$).
    По построению $m_0 > m(\epsilon, k)$.
    А значит $\forall s > m(\epsilon, k) \hookrightarrow$
    \[
        |f_k(x_{n_s}) - f_k(x_0) + f_k(x_{n_{m_0}}) - f_k(x_{n_{m_0}})| < 2\epsilon,
    \]
    что и даёт нужную нам сходимость.

    Осталось показать, что это верно для любого $g \in B_1^*(0) \subset L^*$, то есть что
    \[
        g(x_{n_m}) \rightarrow g(x_0).
    \]
    Пусть это не так.
    Тогда существует подпоследовательность $x_{n_{m_{s}}}$ такая, что для некоторого $\epsilon_0 > 0$
    \[
        |g(x_{n_{m_{s}}}) - g(x_0)| \geq \epsilon_0,
    \]
    то есть $x_{n_{m_{s}}} \notin V(x_0, g, \epsilon_0)$.
    Заметим, что $\{x_{n_{m_{s}}}\}$ -- бесконечное подмножество слабого компакта $K \cap L$ а, следовательно, имеет предельную точку $y_0$.
    Так как $f_k(x_{n_{m_{s}}})$ -- сходящаяся, то, абсолютно аналогично, $f_k(x_{n_{m_{s}}})$ будет сходиться к $f_k(y_0)$.
    По единственности предела $f_k(x_0) = f_k(y_0)$.
    Поэтому $f_k(x_0 - y_0) = 0$.
    А значит в силу плотности $\{f_k\}$ в шаре, то для всякого $f$ из шара $f(x_0 - y_0) = 0$.
    Поэтому $x_0 = y_0$ -- противоречие.

    Таким образом мы получили искомую подпоследовательность.
\end{proof}
\begin{corollary}
    Если $X$ -- рефлексивное пространство, то $B_1(0)$ является слабым секвенциальным компактом.
\end{corollary}
\begin{note}
    Верно и обратно: всякий слабый секвенциальный компакт является и слабым топологическим компактом.
\end{note}